\subsection{\texorpdfstring{THE VECTOR SPACE \(C^{n}\)}{THE VECTOR SPACE C TO THE N}}

Since the topological space \(\mathbf{C}\) can be identified with \(\mathbf{R}^2\) , the topological product \(\mathbf{C}^n\) of \(n\) spaces identical with \(\mathbf{C}\) can be identified with \(\mathbf{R}^{2n}\) quâ topological space; likewise, the topological group structure of \(\mathbf{C}^n\) , the product of the additive (topological) group structures of the \(n\) factors, can be identified with that of the additive group \(\mathbf{R}^{2n}\) . But since \(\mathbf{C}\) is a field, we may define on \(\mathbf{C}^n\) the structure of an \(n\) -dimensional vector space over \(\mathbf{C}\) , the product \(az\) of a complex number \(a\) and a point \(z = (z_i)\) of \(\mathbf{C}^n\) being the point \((az_i)\) ; this vector space structure should be carefully distinguished from the structure of a \(2n\) -dimensional vector space over \(\mathbf{R}\) , defined on \(\mathbf{R}^{2n}\) (Chapter VI, § 1, no. 3). We shall reserve the notation \(\mathbf{C}^n\) for the topological space which is the product of \(n\) spaces identical with \(\mathbf{C}\) , endowed in addition with the vector space structure over \(\mathbf{C}\) which has just been defined; \(\mathbf{C}^n\) is called complex number space of \(n\) dimensions. Note that the mapping \((t, z) \to tz\) is continuous on \(\mathbf{C} \times \mathbf{C}^n\) .

An affine linear mapping of \(C^{n}\) into \(C^{m}\) is also an affine linear mapping of \(R^{2n}\) into \(R^{2m}\) , but the converse is false.

For example, the mapping \(z \rightarrow \overline{z}\) is a linear mapping of the vector space \(R^{2}\) onto itself but is not a linear mapping of the vector space C onto itself.

Every affine linear mapping of \(C^{n}\) into \(C^{m}\) is therefore uniformly continuous; in particular, every affine linear mapping of \(C^{n}\) onto itself is a homeomorphism.

Every affine linear variety of p dimensions \((p \leqslant n)\) in the vector space \(C^{n}\) is also an affine linear variety of 2p dimensions in the vector space \(R^{2n}\) ; here again, the converse is false. To avoid all confusion, (affine) linear varieties of p dimensions in \(C^{n}\) will be called complex linear varieties of p dimensions (the linear varieties of \(R^{2n}\) being called real linear varieties when necessary to prevent misunderstanding). In particular, complex linear varieties of one dimension (resp. two dimensions) are called complex lines (resp. complex planes), and complex linear varieties of n - r dimensions are called complex hyperplanes.

It is often convenient to regard real number space \(R^{n}\) as embedded in complex number space \(C^{n}\) , by identifying \(R^{n}\) with the subset of \(C^{n}\) defined by the relations \(\mathfrak{J}(z_{k}) = o\) ( \(1 \leqslant k \leqslant n\) ). The topological group structure induced on this subset by the topological group structure of \(C^{n}\) coincides with that of \(R^{n}\) .

Note that \(R^{n}\) , thus embedded in \(C^{n}\) , is not a complex linear variety in \(C^{n}\) .

A system of p vectors of \(R^{n}\) which is free over R is also free over C. Every real linear variety V of p dimensions in \(R^{n}\) generates a complex linear variety \(V'\) of p dimensions in \(C^{n}\) such that V is the trace of \(V'\) on \(R^{n}\) ; if V is defined by a system of n-p linear equations \(f_{k}(x)=a_{k}\) , where the \(f_{k}\) are linear forms on \(R^{n}\) (with real coefficients, and linearly independent over R) and the \(a_{k}\) are real numbers, then the same equations define \(V'\) , but now the coordinates of x take complex values.

Conversely, if a complex linear variety of \(p\) dimensions has a non-empty intersection with \(\mathbf{R}^n\) , this intersection is a real linear variety, but its dimension may be \(< p\) .

